\documentclass[11pt]{amsart}

\usepackage{amsmath,amssymb,amsthm,mathtools}
\usepackage{mathrsfs}
\usepackage{hyperref}

\DeclareMathOperator{\Sp}{Sp}
\DeclareMathOperator{\Stab}{Stab}

\newtheorem{theorem}{Theorem}[section]
\newtheorem{proposition}[theorem]{Proposition}
\newtheorem{corollary}[theorem]{Corollary}
\theoremstyle{remark}
\newtheorem{remark}[theorem]{Remark}

\title{Minimal-Good-Lattice Normalization of Lattice-Model Splittings}
\author{Author Placeholder}
\subjclass[2020]{22E50, 22E45}
\keywords{metaplectic extension, reductive dual pair, lattice model, building}
\date{}

\begin{document}

\begin{abstract}
We compare the splittings of the universal metaplectic extension furnished by
generalized lattice models. Their ratios are continuous
$\mathbb C^\times$-valued characters satisfying a cocycle identity. A minimal
good lattice on the opposite member of the dual pair then defines a
choice-independent splitting, and comparison with it gives the normalized
character $\xi_{x,x'}$.
\end{abstract}

\maketitle

\section{Setup and known inputs}

Let
\[
(G,G')=(U(V,h),U(V',h'))\subseteq \Sp(W)
\]
be the type-I reductive dual pair under consideration, and work in the
universal extension
\[
1\longrightarrow \mathbb C^\times
\longrightarrow \widetilde{\Sp}(W)
\longrightarrow \Sp(W)
\longrightarrow 1.
\]
Fix $x\in\mathcal B(G)$ and put $K_x:=\Stab_G(x)$. We assume that the
building has been refined so that $K_x$ is constant on each open facet. For
every $x'\in\mathcal B(G')$, the generalized lattice model gives a continuous
group-homomorphic splitting
\[
s_{x,x'}\colon K_x\longrightarrow\widetilde{\Sp}(W).
\]

We use the following known inputs. The nonempty set $\mathscr M(V')$ of
minimal good lattices is a single $G'$-orbit. Each
$\Lambda'\in\mathscr M(V')$ determines a normalized building point
$x'_{\Lambda'}$, equivariantly in the sense that
\[
x'_{a\Lambda'}=a x'_{\Lambda'}\qquad(a\in G').
\]
Generalized-lattice-model naturality says that, for a lift $\widetilde a$ of
$1\otimes a$ and $g\in K_x$,
\begin{equation}
s_{x,ax'}(g)=\widetilde a\,s_{x,x'}(g)\widetilde a^{-1}.
\label{eq:naturality}
\end{equation}
Finally, the known commuting-cover theorem for this reductive dual pair says
that the pullback covering groups over $G$ and $G'$ commute elementwise in
$\widetilde{\Sp}(W)$. This nontrivial cover-level theorem is used as an input;
it is not deduced from commutativity in $\Sp(W)$.

\section{Comparison characters}

For $x'_1,x'_2\in\mathcal B(G')$ define
\begin{equation}
\xi_{x;x'_1,x'_2}(g)
:=s_{x,x'_1}(g)\bigl(s_{x,x'_2}(g)\bigr)^{-1},
\qquad g\in K_x.
\label{eq:raw-xi}
\end{equation}
The two factors have the same image $g$ in $\Sp(W)$, so their quotient lies
in the central kernel $\mathbb C^\times$.

\begin{proposition}
The map $\xi_{x;x'_1,x'_2}\colon K_x\to\mathbb C^\times$ is a continuous
character. For all $x'_1,x'_2,x'_3\in\mathcal B(G')$ it satisfies
\[
\xi_{x;x'_1,x'_3}
=\xi_{x;x'_1,x'_2}\,\xi_{x;x'_2,x'_3}.
\]
\end{proposition}

\begin{proof}
Write $s_i=s_{x,x'_i}$ and $\xi_{ij}=\xi_{x;x'_i,x'_j}$. Since
$s_i(g)=\xi_{ij}(g)s_j(g)$ and $\xi_{ij}(g)$ is central, multiplicativity of
the splittings gives
\[
\begin{aligned}
\xi_{ij}(gh)s_j(g)s_j(h)
 &=s_i(gh)=s_i(g)s_i(h)\\
 &=\xi_{ij}(g)\xi_{ij}(h)s_j(g)s_j(h).
\end{aligned}
\]
Cancellation proves $\xi_{ij}(gh)=\xi_{ij}(g)\xi_{ij}(h)$; continuity follows
from continuity of the two splittings. Direct cancellation in the central
kernel gives
\[
s_1(g)s_3(g)^{-1}
=\bigl(s_1(g)s_2(g)^{-1}\bigr)
 \bigl(s_2(g)s_3(g)^{-1}\bigr),
\]
which is the cocycle identity.
\end{proof}

\section{The minimal splitting}

For $\Lambda'\in\mathscr M(V')$ put
\[
s^{\min}_{x,\Lambda'}:=s_{x,x'_{\Lambda'}}.
\]

\begin{theorem}[Independence of the minimal good lattice]
For all $\Lambda'_0,\Lambda'_1\in\mathscr M(V')$,
\[
s^{\min}_{x,\Lambda'_0}=s^{\min}_{x,\Lambda'_1}
\]
as continuous group-homomorphic splittings. Their common value therefore
defines a canonical splitting
\[
s_x^{\min}\colon K_x\longrightarrow\widetilde{\Sp}(W).
\]
\end{theorem}

\begin{proof}
Choose $a\in G'$ with $\Lambda'_1=a\Lambda'_0$, and put
$x'_0=x'_{\Lambda'_0}$. Equivariance gives
$x'_{\Lambda'_1}=a x'_0$. If $\widetilde a$ is any lift of $1\otimes a$,
then lattice-model naturality \eqref{eq:naturality} gives
\begin{equation}
s_{x,a x'_0}(g)
=\widetilde a\,s_{x,x'_0}(g)\widetilde a^{-1}.
\label{eq:conjugacy}
\end{equation}
This is only a conjugacy statement. The element $s_{x,x'_0}(g)$ lies in the
pullback cover over $G$, while $\widetilde a$ lies in the pullback cover over
$G'$. The separate commuting-cover theorem therefore gives
\[
[\widetilde a,s_{x,x'_0}(g)]=1.
\]
Substitution in \eqref{eq:conjugacy} yields
$s_{x,a x'_0}(g)=s_{x,x'_0}(g)$ for every $g\in K_x$. Since
$x'_{\Lambda'_1}=a x'_0$, the two minimal splittings agree. Their common
value remains continuous, group-homomorphic, and a splitting because every
$s_{x,x'}$ has these properties by construction.
\end{proof}

\section{Normalized correction character}

\begin{corollary}
For arbitrary $x'\in\mathcal B(G')$, the formula
\begin{equation}
\boxed{\quad
\xi_{x,x'}(g)
:=s_x^{\min}(g)\bigl(s_{x,x'}(g)\bigr)^{-1},
\qquad g\in K_x,
\quad}
\label{eq:normalized-xi}
\end{equation}
defines a continuous $\mathbb C^\times$-valued character, and
\begin{equation}
\boxed{\quad s_x^{\min}=\xi_{x,x'}s_{x,x'}.\quad}
\label{eq:normalization}
\end{equation}
Moreover,
\begin{equation}
\xi_{x;x'_1,x'_2}
=\xi_{x,x'_1}^{-1}\xi_{x,x'_2},
\label{eq:two-xi}
\end{equation}
and $\xi_{x,x'_{\Lambda'}}=1$ for every minimal good lattice $\Lambda'$.
\end{corollary}

\begin{proof}
Equation \eqref{eq:normalized-xi} is the comparison of the two splittings
$s_x^{\min}$ and $s_{x,x'}$, so the proposition proves that it is a continuous
character. Rearranging its definition gives \eqref{eq:normalization}. Also,
centrality gives
\[
\begin{aligned}
\xi_{x,x'_1}^{-1}(g)\xi_{x,x'_2}(g)
 &=s_{x,x'_1}(g)(s_x^{\min}(g))^{-1}
   s_x^{\min}(g)(s_{x,x'_2}(g))^{-1}\\
 &=s_{x,x'_1}(g)(s_{x,x'_2}(g))^{-1}
  =\xi_{x;x'_1,x'_2}(g),
\end{aligned}
\]
which proves \eqref{eq:two-xi}. If $x'=x'_{\Lambda'}$, then independence of
$\Lambda'$ gives $s_{x,x'}=s_x^{\min}$, hence $\xi_{x,x'}=1$.
\end{proof}

\begin{remark}[Scope of the normalization]
The formal comparison above does not prove that $\xi_{x,x'}$ is trivial on
$K_x^+$, factors through $K_x/K_x^+$, is quadratic, or equals a determinant
product on a finite graded quotient. Each such refinement requires a separate
argument.
\end{remark}

\end{document}
